\documentclass{ximera}
\title{Point-Slope}

\newcommand{\pskip}{\vskip 0.1 in}

\begin{document}
\begin{abstract}
The point-slope equation of a line from the perspective of linear functions.
\end{abstract}
\maketitle

\begin{example} \label{Eldfm3fg3}
Suppose $y=f(x)$ is a linear function with constant rate of change
\[
  \frac{\Delta y}{\Delta x} = \frac{2}{3}.
\]
Suppose also that $f(4) = 7$.

\begin{enumerate}
\item Evaluate $f(23)$.

\item Find an expression for $f(x)$.

\end{enumerate}

\begin{explanation}
\begin{enumerate}
\item Since the average rate of change of $f(x)$ between the inputs $x=4$ and $x=23$ is
\[
     \frac{\Delta y}{\Delta x}  = \frac{f(23) - f(4)}{23-4} = \frac{2}{3} ,
\]
\[
        f(23) - f(4) = \frac{2}{3}\left( 23-4 \right),
\]
and
\begin{align*}
       f(23) &= f(4) + \frac{2}{3}\left( 23-4 \right) \\
               &= \answer{7} + \frac{2}{3}\left( 19 \right) \\
               &= 7 + \answer{\frac{38}{3}} \\
               & = \frac{\answer{21}}{3} + \frac{38}{3} \\
               & = \answer{\frac{59}{3}} .
\end{align*}

\item We follow the idea in part (a) to find an expression for $f(x)$. 
Since the average rate of change of $f(x)$ between the inputs $4$ and $x$ ($x\neq 4$) is
\[
     \frac{\Delta y}{\Delta x}  = \frac{f(x) - f(4)}{x-4} = \frac{2}{3} ,
\]
\[
        f(x) - f(4) = \frac{2}{3}\left( x-4 \right),
\]
and
\[
  f(x) = f(4) + \frac{2}{3}\left( x-4 \right).
\]
Note this last equation also holds if $x=4$. So this is an expression for $f(x)$.
\end{enumerate}
\end{explanation}
\end{example}


\begin{example} \label{Exfef3fr3m}
Suppose $y=f(x)$ is a linear function with constant rate of change
\[
  \frac{\Delta y}{\Delta x}.
\]
Then for any number $a$ in the domain of $f$, the same reasoning as in Example 1 shows
\[ 
    f(x) = f(a) + \frac{\Delta y}{\Delta x} \left( x-a \right).
\]
This is the \emph{point-slope} expression for the linear function $y=f(x)$. It says to evaluate $f(x)$ we start with $f(a)$ and add to that output the change
\[
 \Delta y =  f(x) - f(a) = \frac{\Delta y}{\Delta x} \left( x-a \right)
\]
in the output as the input changes from $a$ to $x$. This change is equal to the product of the constant rate of change $\Delta y/\Delta x$ of $y$ with respect to $x$ and the change
\[
 \Delta x = x-a
\]
in $x$.
\end{example}

\begin{example} \label{Ex3er3}
The function $f$ is linear with $f(-2) = 8$ and $f(5) = -3$. Find an expression for $f(x)$.

\begin{explanation}
Since the function $y=f(x)$ is linear, the rate of change of $y$ with respect to $x$ is constant. This rate of change is
\begin{align*}
    \frac{\Delta y}{\Delta x} &= \frac{f(5) - f(-2)}{5 - (-2)} \\
                                        &= \frac{-3 - 8}{7} \\
                                        &= -\frac{11}{7}.
\end{align*}
So
\begin{align*}
               f(x) &= f(-2) + \frac{\Delta y}{\Delta x}\left( x - (-2) \right) \\
                     &= 8 - \frac{11}{7} \left( x+2 \right).
\end{align*}

\emph{Note:} We could have also based our expression for $f(x)$ using $f(5) = -3$.

Find such an expression and then use algebra to show this expression for $f(x)$ is equivalent to the one above. 

\begin{freeResponse}
\end{freeResponse}
\end{explanation}
\end{example}

\end{document}